\section*{Chapitre \ref{ch:g7:atoms-and-molecules} --- Atomes et molécules}

\begin{solution}{exo:g7:atoms-and-molecules:1}
Un \omterm{def:g7:atoms-and-molecules:atom}{atome} est l'une des particules extrêmement petites à partir
desquelles toute la matière est construite (il en existe une centaine
de sortes). Une \omterm{def:g7:atoms-and-molecules:molecule}{molécule} est un groupe d'\omterm{def:g7:atoms-and-molecules:atom}{atomes} solidement liés entre
eux, le plus petit morceau d'une substance qui soit encore cette
substance.
\end{solution}

\begin{solution}{exo:g7:atoms-and-molecules:2}
\ce{H}, \ce{C}, \ce{N}, \ce{O}, \ce{Na}, \ce{Fe}.
\end{solution}

\begin{solution}{exo:g7:atoms-and-molecules:3}
1 \omterm{def:g7:atoms-and-molecules:atom}{atome} de carbone et 2 \omterm{def:g7:atoms-and-molecules:atom}{atomes} d'oxygène~: 3 \omterm{def:g7:atoms-and-molecules:atom}{atomes} en tout.
\end{solution}

\begin{solution}{exo:g7:atoms-and-molecules:4}
Eau, dioxygène, méthane, diazote.
\end{solution}

\begin{solution}{exo:g7:atoms-and-molecules:5}
\ce{NO2}.
\end{solution}

\begin{solution}{exo:g7:atoms-and-molecules:6}
\ce{O2} est une \omterm{def:g7:atoms-and-molecules:molecule}{molécule} formée de deux \omterm{def:g7:atoms-and-molecules:atom}{atomes} d'oxygène liés~; \ce{2O}
désigne deux \omterm{def:g7:atoms-and-molecules:atom}{atomes} d'oxygène séparés. \ce{CO} est une \omterm{def:g7:atoms-and-molecules:molecule}{molécule} formée
d'un \omterm{def:g7:atoms-and-molecules:atom}{atome} de carbone et d'un \omterm{def:g7:atoms-and-molecules:atom}{atome} d'oxygène (deux majuscules, deux
\omterm{def:g7:atoms-and-molecules:symbol}{symboles})~; \ce{Co} est un \omterm{def:g7:atoms-and-molecules:atom}{atome} de cobalt (un seul \omterm{def:g7:atoms-and-molecules:symbol}{symbole}, une
majuscule puis une minuscule).
\end{solution}

\begin{solution}{exo:g7:atoms-and-molecules:7}
\ce{3H2O}~: $3\times 2 = 6$ \omterm{def:g7:atoms-and-molecules:atom}{atomes} d'hydrogène et $3\times 1 = 3$ \omterm{def:g7:atoms-and-molecules:atom}{atomes}
d'oxygène. \ce{5O2}~: pas d'hydrogène, $5\times 2 = 10$ \omterm{def:g7:atoms-and-molecules:atom}{atomes} d'oxygène.
\end{solution}

\begin{solution}{exo:g7:atoms-and-molecules:8}
Le noir est le carbone, le rouge l'oxygène~: un carbone, deux oxygènes,
\ce{CO2}, le dioxyde de carbone.
\end{solution}

\begin{solution}{exo:g7:atoms-and-molecules:9}
L'eau de Dalton était \ce{HO}. La vraie \omterm{def:g7:atoms-and-molecules:molecule}{molécule}, \ce{H2O}, a un \omterm{def:g7:atoms-and-molecules:atom}{atome}
d'hydrogène de plus~: il manquait un \omterm{def:g7:atoms-and-molecules:atom}{atome} à la \omterm{def:g7:atoms-and-molecules:molecule}{molécule} de Dalton.
\end{solution}

\begin{solution}{exo:g7:atoms-and-molecules:10}
L'eau pure n'est pas un \omterm{def:g6:pure-substances-and-mixtures:mixture}{mélange}~: toutes ses \omterm{def:g7:atoms-and-molecules:molecule}{molécules} sont des
\omterm{def:g7:atoms-and-molecules:molecule}{molécules} \ce{H2O} identiques. L'\omterm{def:g4:air-a-mixture-of-gases:air}{air} est un \omterm{def:g6:pure-substances-and-mixtures:mixture}{mélange}~: il contient des
\omterm{def:g7:atoms-and-molecules:molecule}{molécules} de diazote, des \omterm{def:g7:atoms-and-molecules:molecule}{molécules} de dioxygène, des \omterm{def:g7:atoms-and-molecules:atom}{atomes} d'argon et
quelques \omterm{def:g7:atoms-and-molecules:molecule}{molécules} de dioxyde de carbone.
\end{solution}

\begin{solution}{exo:g7:atoms-and-molecules:11}
Dix millions d'\omterm{def:g7:atoms-and-molecules:atom}{atomes} font \qty{1}{mm}, donc \qty{1}{cm} $= \qty{10}{mm}$
en demande $10 \times$ dix millions $=$ cent millions d'\omterm{def:g7:atoms-and-molecules:atom}{atomes}.
$\qty{100}{m} = \num{100000}$~mm en demande $\num{100000}\times$ dix
millions $=$ un million de millions d'\omterm{def:g7:atoms-and-molecules:atom}{atomes}.
\end{solution}

\begin{solution}{exo:g7:atoms-and-molecules:12}
$6 + 12 + 6 = 24$ \omterm{def:g7:atoms-and-molecules:atom}{atomes}. Hydrogène~: $\frac{12}{24} = \qty{50}{\percent}$.
Carbone~: $\frac{6}{24} = \qty{25}{\percent}$ (et oxygène les
\qty{25}{\percent} restants).
\end{solution}

\begin{solution}{pb:g7:atoms-and-molecules:1}
\textbf{1.} Diazote \ce{N2}, 2 \omterm{def:g7:atoms-and-molecules:atom}{atomes}~; dioxygène \ce{O2}, 2 \omterm{def:g7:atoms-and-molecules:atom}{atomes}~;
dioxyde de carbone \ce{CO2}, 3 \omterm{def:g7:atoms-and-molecules:atom}{atomes}.

\textbf{2.} Les \omterm{def:g7:atoms-and-molecules:atom}{atomes} d'argon ne s'assemblent pas en \omterm{def:g7:atoms-and-molecules:molecule}{molécules}~:
l'argon est fait d'\omterm{def:g7:atoms-and-molecules:atom}{atomes} isolés.

\textbf{3.} Azote~: $78 \times 2 = 156$ \omterm{def:g7:atoms-and-molecules:atom}{atomes}. Oxygène~: $21\times 2 =
42$ \omterm{def:g7:atoms-and-molecules:atom}{atomes}.

\textbf{4.} $156 + 42 + 1 = 199$ \omterm{def:g7:atoms-and-molecules:atom}{atomes}.

\textbf{5.} $\frac{156}{199} \approx 0.78$, soit environ
\qty{78}{\percent}.

\textbf{6.} Dans les \omterm{def:g7:atoms-and-molecules:molecule}{molécules} de dioxygène~: $16\times 2 = 32$~; dans
les \omterm{def:g7:atoms-and-molecules:molecule}{molécules} de dioxyde de carbone~: $4 \times 2 = 8$~; en tout
$32 + 8 = 40$.

\textbf{7.} Il manque $42 - 40 = 2$ \omterm{def:g7:atoms-and-molecules:atom}{atomes} d'oxygène~: l'équivalent
d'une \omterm{def:g7:atoms-and-molecules:molecule}{molécule} de dioxygène. Ils sont partis dans des \omterm{def:g7:atoms-and-molecules:molecule}{molécules} d'eau,
\ce{H2O}, que le corps fabrique avec l'hydrogène de ses aliments~; cette
eau sort en partie sous forme de vapeur d'eau dans le souffle (retirée
dans ces chiffres).

\textbf{8.} \omterm{def:g4:air-a-mixture-of-gases:air}{Air} inspiré~: 0. \omterm{def:g4:air-a-mixture-of-gases:air}{Air} expiré~: 4 (un dans chaque \ce{CO2}).

\textbf{9.} Des aliments que le corps utilise~: les sucres et les
graisses contiennent des \omterm{def:g7:atoms-and-molecules:atom}{atomes} de carbone, et le corps les rejette dans
le dioxyde de carbone qu'il expire.

\textbf{10.} $21 - 16 = 5$ \omterm{def:g7:atoms-and-molecules:molecule}{molécules} de dioxygène ont disparu et 4
\omterm{def:g7:atoms-and-molecules:molecule}{molécules} de dioxyde de carbone sont apparues.

\textbf{11.} Diazote~: deux boules bleues reliées par trois bâtonnets.
Dioxygène~: deux boules rouges reliées par deux bâtonnets. Dioxyde de
carbone~: une boule noire entre deux boules rouges, chacune reliée à
elle par deux bâtonnets, en ligne droite. Argon~: une seule boule bleu
pâle.

\textbf{12.} $156 + 32 + 12 + 1 = 201$ \omterm{def:g7:atoms-and-molecules:atom}{atomes} (les 4 \omterm{def:g7:atoms-and-molecules:molecule}{molécules} de
dioxyde de carbone contiennent 12 \omterm{def:g7:atoms-and-molecules:atom}{atomes}). C'est $201 - 199 = 2$ \omterm{def:g7:atoms-and-molecules:atom}{atomes}
de plus que l'\omterm{def:g4:air-a-mixture-of-gases:air}{air} inspiré~: 4 \omterm{def:g7:atoms-and-molecules:atom}{atomes} de carbone gagnés grâce aux
aliments, moins les 2 \omterm{def:g7:atoms-and-molecules:atom}{atomes} d'oxygène partis dans l'eau.
\end{solution}
